\subsection{Sprint 4}

Na sprint final, procuramos evitar a repetição das mobilizações de última hora
que haviam marcado as entregas anteriores. Nós, alunos AGES IV, cobramos que as
funcionalidades planejadas fossem concluídas entre dez dias e duas semanas antes
da apresentação final às \textit{stakeholders}. Essa margem foi definida para que
o encerramento do desenvolvimento não coincidisse com a preparação da entrega e
para que houvesse tempo de testar a aplicação integrada, verificar pendências e
organizar uma demonstração coerente do produto.

A disponibilização antecipada das histórias das Sprints 3 e 4, somada ao
acompanhamento mais frequente adotado depois da Sprint 2, contribuiu para que o
time concluísse o desenvolvimento antes do prazo interno estabelecido. Com isso,
a etapa final não ficou concentrada em integrar grandes quantidades de código na
véspera. Pela primeira vez no semestre, conseguimos reservar um período
significativo para avaliar o sistema já reunido e disponível no ambiente de
produção.

Minha principal atividade nessa etapa foi testar a aplicação em produção. Eu
percorri os fluxos disponíveis, comparei o comportamento encontrado com o que
esperávamos apresentar e procurei identificar funcionalidades ausentes e erros.
Esses testes revelaram algumas pendências de resolução rápida, que puderam ser
encaminhadas antes da entrega. Além de validar o produto, esse processo ajudou a
definir quais jornadas estavam estáveis e deveriam compor a demonstração para as
\textit{stakeholders}.

Na reta final, concentramo-nos na preparação da apresentação e na organização do
fluxo da demonstração. Como o desenvolvimento havia sido encerrado com
antecedência, os demais integrantes também tiveram tempo para finalizar seus
relatórios sem competir diretamente com correções urgentes no produto. Esse
cenário foi diferente das sprints anteriores, nas quais a preparação da
apresentação ocorria ao mesmo tempo que a revisão de PRs e a tentativa de concluir
tarefas abertas.

A apresentação final transcorreu conforme o planejado, e as \textit{stakeholders}
ficaram muito satisfeitas com o resultado entregue. Além do retorno positivo
sobre o projeto, elas compartilharam sugestões sobre como apresentar melhor a
proposta e construir um \textit{pitch} mais convincente para a tentativa de
alcançar o reconhecimento de projeto destaque. Após o encerramento da entrega,
também dediquei tempo à elaboração do relatório final, registrando as atividades
e os resultados do semestre.

\medskip

\fbox{%
    \parbox{0.95\linewidth}{%
        \setlength{\parindent}{15pt}
        \itshape A Sprint 4 foi consideravelmente mais tranquila porque, desta
        vez, trabalhamos com uma margem anterior à data da apresentação. Concluir
        o desenvolvimento mais cedo permitiu testar o sistema em produção,
        corrigir pendências menores e ensaiar uma demonstração que representasse
        adequadamente o que havia sido construído. A diferença em relação às
        entregas anteriores mostrou que parte da pressão que enfrentamos não era
        inevitável, mas consequência de integrar e validar o trabalho tarde
        demais.

        Considero que nós, AGES IV, conseguimos aplicar aprendizados das sprints
        anteriores. Passamos a cobrar as entregas com antecedência, acompanhar os
        PRs mais de perto e reservar tempo explícito para testes e preparação da
        apresentação. O resultado não apaga os problemas de organização que
        ocorreram durante o semestre, especialmente no ClickUp, mas demonstra que
        o time conseguiu ajustar sua forma de trabalho e encerrar o projeto com
        um processo mais controlado.

        Para mim, o principal aprendizado dessa sprint foi que o papel de gestão
        não termina quando as tarefas de desenvolvimento são concluídas. A
        qualidade percebida na entrega também depende da validação do produto, da
        escolha de um fluxo de demonstração confiável e da capacidade de explicar
        seu valor às \textit{stakeholders}. O retorno recebido na apresentação e
        as orientações para o \textit{pitch} reforçaram que uma boa solução
        técnica precisa ser acompanhada por uma comunicação igualmente bem
        preparada.
    }%
}
